% =====================================================================
%  ch02a.tex —— 第 2 章标题 + §2.1（原书 p.530–538）
%  底本：Orlov, Derived categories of coherent sheaves and equivalences
%        between them, Russian Math. Surveys 58:3 (2003), 511–591
%  本块止于 “2.2. Examples of equivalences: flopping birational
%  transformations” 标题之前（2.2 归 ch02b）。
% =====================================================================
% [原文 p.530]
\chapter{凝聚层范畴及其间函子}

\section{凝聚层范畴的基本性质}
\takeaway{本节把 Serre 对偶 (6)、由积 $X\times Y$ 上的对象给出的函子 $\Phi_{\cE}$ 与 $\Psi_{\cE}$（连同伴随 (8)、复合律 (9)）整理成一套函子演算；再由对角线上典范层之幂造出双分次代数 $\mathrm{HA}(X)$，并证明它、霍奇菱形的竖列、陈特征与 $K$ 群都是导出范畴的不变量。}

从现在起，我们只考虑光滑完备代数簇上凝聚层的有界导出范畴。为简省起见，我们总把 $\Db(\coh(X))$ 简写作 $\Db(X)$。此外，若函子正合——例如平坦态射下的逆像，或与局部自由层作张量积——我们就省略导出函子的符号。

对维数为 $n$ 的光滑完备簇 $X$，凝聚层的有界导出范畴带有 Serre 函子（见定义 1.1.3），它由 $(\,\cdot\,)\otimes\omega_X[n]$ 给出，其中 $\omega_X$ 是典范层（见 [6]）。于是，对任意一对对象 $\cE,\cF\in\Db(X)$，有同构
\begin{equation}
\Hom(\cE,\cF)=\Hom(\cF,\cE\otimes\omega_X[n])^{*} \tag{6}
\end{equation}

\yizhu{Serre 函子指三角范畴 $\cD$ 上一个正合自函子 $S$，连同对任意 $A,B$ 自然的同构 $\Hom(A,B)\cong\Hom(B,S(A))^{*}$；它等价于说 $\cD$ 上存在 Serre 对偶。对光滑完备簇 $X$，$\Db(X)$ 的 Serre 函子就是 $(\,\cdot\,)\otimes\omega_X[n]$，$n=\dim X$——这正是 Serre 对偶定理在导出范畴里的整体形式，式 (6) 只是它的逐字展开。下文 (13) 处的论证（等价比与 Serre 函子交换）与本节末的 (16) 都依赖 $S$ 的存在性。〔来源：原书 §2.1 式 (6) 与定义 1.1.3；Hartshorne, \emph{Algebraic Geometry}, III.7（A 级）〕}

如上一节所表明的，光滑完备代数簇之间的每个态射 $f:X\to Y$ 都诱导两个正合函子，即直接像函子
% [原文 p.531]
$\Rf f_*:\Db(X)\to\Db(Y)$ 与逆像函子 $\Lf f^*:\Db(Y)\to\Db(X)$，且这两个函子互为伴随。此外，每个对象 $\cE\in\Db(X)$ 都定义正合张量积函子 $\otL\cE:\Db(X)\to\Db(X)$。

我们可以用这些标准的导出函子，在导出范畴 $\Db(X)$ 与 $\Db(Y)$ 之间引入一大类新的正合函子。

设 $X$ 与 $Y$ 是域 $k$ 上两个光滑完备簇，维数分别为 $n$ 与 $m$。考虑积 $X\times Y$，并记 $p$ 与 $q$ 为 $X\times Y$ 到 $X$ 与到 $Y$ 的投影：
\begin{equation*}
X\xleftarrow{\ p\ }X\times Y\xrightarrow{\ q\ }Y.
\end{equation*}
每个对象 $\cE\in\Db(X\times Y)$ 都决定一个从导出范畴 $\Db(X)$ 到导出范畴 $\Db(Y)$ 的正合函子 $\Phi_{\cE}$，它由公式
\begin{equation}
\Phi_{\cE}(\,\cdot\,):=\Rb q_*\bigl(\cE\otL p^*(\,\cdot\,)\bigr) \tag{7}
\end{equation}
给出。此外，对同一个对象 $\cE\in\Db(X\times Y)$，还可以配上另一个从导出范畴 $\Db(Y)$ 到导出范畴 $\Db(X)$ 的函子 $\Psi_{\cE}$，它由与 (7) 类似的规则
\begin{equation*}
\Psi_{\cE}(\,\cdot\,):=\Rf p_*\bigl(\cE\otL q^*(\,\cdot\,)\bigr)
\end{equation*}
定义。

容易验证，函子 $\Phi_{\cE}$ 既有左伴随，也有右伴随。

\yizhu{称 $\cE$ 为函子 $\Phi_{\cE}$ 的\emph{核}（kernel object），于是 $\Phi_{\cE}$ 就是积分变换：先用 $p^{*}$ 把 $X$ 上的对象拉到积上，再与核作导出张量积，最后沿 $q$ 推回 $Y$。$\Psi_{\cE}$ 是同一核在相反方向上的变换，故 $\Phi$ 与 $\Psi$ 只在 $p$、$q$ 的次序上不同。把核取成 $\cO_{\Delta}$ 得到恒等函子；取成 Poincaré 丛则得到 Fourier--Mukai 变换——参见引言与第 5 章。〔来源：原书 §2.1 式 (7) 及其后一段；(7) 在引言式 (1) 处已出现，本处补出核的说法〕}

\begin{lemma}{2.1.1}
函子 $\Phi_{\cE}$ 有左伴随函子 $\Phi^{*}_{\cE}$ 与右伴随函子 $\Phi^{!}_{\cE}$，它们分别由公式
\begin{equation}
\Phi^{*}_{\cE}\cong\Psi_{\cE^{\vee}\otimes q^{*}\omega_Y[m]}
\qquad\text{和}\qquad
\Phi^{!}_{\cE}\cong\Psi_{\cE^{\vee}\otimes p^{*}\omega_X[n]} \tag{8}
\end{equation}
给出。
\end{lemma}

这里 $\omega_X$ 与 $\omega_Y$ 分别是 $X$ 与 $Y$ 上的典范层，而 $\cE^{\vee}$ 是 $\RHom(\cE,\cO_{X\times Y})$ 的简便记法。

\begin{proof}
我们给出左伴随函子的证明。它来自下面一串同构：
\begin{equation*}
\begin{aligned}
\Hom\bigl(A,\Rf q_*(\cE\otL p^{*}B)\bigr)
&\cong\Hom(q^{*}A,\cE\otL p^{*}B)\\
&\cong\Hom\bigl(p^{*}B,\cE^{\vee}\otL q^{*}A\otimes\omega_{X\times Y}[n+m]\bigr)^{*}\\
&\cong\Hom\bigl(B,\Rf p_*(\cE^{\vee}\otL q^{*}(A\otimes\omega_Y[m]))\otimes\omega_X[n]\bigr)^{*}\\
&\cong\Hom\bigl(\Rf p_*(\cE^{\vee}\otL q^{*}(A\otimes\omega_Y[m])),B\bigr).
\end{aligned}
\end{equation*}
这里用到了直接像与逆像函子之间的伴随、Serre 对偶 (6)（用了两次），以及公式 (5)。
\end{proof}

当然，我们要指出，形如
\begin{equation*}
X\xleftarrow{\ p\ }Z\xrightarrow{\ q\ }Y
\end{equation*}
的任一图表，以及任一对象 $\cE\in\Db(Z)$，都可以配上从 $X$ 上凝聚层的导出范畴到 $Y$ 上凝聚层的导出范畴的函子，它由与 (7) 类似的公式给出。然而，任何这样的函子都同构于形如 (7) 的函子，
% [原文 p.532]
只是把其中的对象换成 $X\times Y$ 上的 $\Rf(p,q)_*\cE$，这里 $(p,q)$ 是由 $Z$ 到直积 $X\times Y$ 的典范态射。

现在设 $X$、$Y$、$Z$ 是三个光滑完备簇，而 $\cE$、$\cF$、$\cG$ 分别是导出范畴 $\Db(X\times Y)$、$\Db(Y\times Z)$、$\Db(X\times Z)$ 中的对象。考虑如下的投影图：
\begin{center}
\begin{tikzcd}[row sep=3.6em,column sep=2.4em,ampersand replacement=\&]
\& X\times Y\times Z \arrow[ld, "p_{12}"] \arrow[d, "p_{13}"] \arrow[rd, "p_{23}"] \& \\
X\times Y \arrow[d, "p^{1}_{12}"] \arrow[rd, "p^{2}_{12}"] \& X\times Z \arrow[ld, "p^{1}_{13}"] \arrow[rd, "p^{3}_{13}"] \& Y\times Z \arrow[ld, "p^{2}_{23}"] \arrow[d, "p^{3}_{23}"] \\
X \& Y \& Z
\end{tikzcd}
\end{center}

对象 $\cE$、$\cF$、$\cG$ 定义三个函子
\begin{equation*}
\Phi_{\cE}:\Db(X)\to\Db(Y),\qquad
\Phi_{\cF}:\Db(Y)\to\Db(Z),\qquad
\Phi_{\cG}:\Db(X)\to\Db(Z),
\end{equation*}
它们由公式 (7) 给出，即
\begin{equation*}
\begin{gathered}
\Phi_{\cE}:=\Rf\pi^{2}_{12*}\bigl(\cE\otL\pi^{1*}_{12}(\,\cdot\,)\bigr),
\qquad
\Phi_{\cF}:=\Rf\pi^{3}_{23*}\bigl(\cF\otL\pi^{2*}_{23}(\,\cdot\,)\bigr)\\
\text{以及}\qquad
\Phi_{\cG}:=\Rf\pi^{3}_{13*}\bigl(\cG\otL\pi^{1*}_{13}(\,\cdot\,)\bigr).
\end{gathered}
\end{equation*}

我们考虑对象 $p^{*}_{12}\cE\otL p^{*}_{23}\cF\in\Db(X\times Y\times Z)$，下文总把它记作 $\cE\boxtimes_Y\cF$。下述断言给出由积上对象代表的导出范畴之间正合函子的复合律。

\yizhu{$\cE\boxtimes_Y\cF$ 不是新的构造，按本处定义它就是沿投影拉回到三重积上再作导出张量积：$\cE\boxtimes_Y\cF:=p^{*}_{12}\cE\otL p^{*}_{23}\cF$。指标的 $Y$ 只记录“两个对象在中间的 $Y$ 方向上被粘合”这一事实。同一个符号 $\boxtimes$ 在本节后文（p.534 起）用于两个\emph{互不相干}的因子之间的外张量积，那时没有下标，含义是 $p^{*}_{13}(\cE_1)\otL p^{*}_{24}(\cE_2)$。〔来源：原书 §2.1 式 (9) 前后的定义〕}

\begin{proposition}{2.1.2}
函子的复合 $\Phi_{\cF}\circ\Phi_{\cE}$ 同构于由
\begin{equation}
\cG=\Rf p_{13*}\bigl(\cE\boxtimes_Y\cF\bigr) \tag{9}
\end{equation}
代表的函子 $\Phi_{\cG}$。
\end{proposition}

证明是直接验证。

% [Tutor Note: 式 (9) 页图印作 $\mathbf{R}p_{13*}$（小写 $p$），而同一段落中三个函子印作 $\mathbf{R}\pi^{3}_{13*}$ 等（希腊 $\pi$，带上标）；文字层同样抽成 $Rp_{13*}$。这是原书记号不统一，本译本照录页图，不改成 $\pi$。]

于是，每个光滑完备代数簇都配上其凝聚层的导出范畴；而两个簇的积上的每个对象 $\cE\in\Db(X\times Y)$ 都配上从三角范畴 $\Db(X)$ 到三角范畴 $\Db(Y)$ 的正合函子 $\Phi_{\cE}$，复合律即上述刚刚描述的那一条。

下述问题对理解这一对应是根本的：
\begin{enumerate}[label=\arabic*),leftmargin=2.4em,itemsep=1pt,topsep=2pt]
\item 两个不同的光滑完备代数簇上凝聚层的导出范畴，何时作为三角范畴等价？
\item 给定簇 $X$ 上凝聚层导出范畴的正合自等价群是什么？（这里指的是正合自等价的同构类构成的群。）
\item 凝聚层导出范畴之间的每个正合函子都可由积上的对象代表，即都具有 (7) 的形式吗？
\end{enumerate}

这一方向上已有一些结果。例如，当簇具有丰富的典范层或反典范层时，前两个问题都能得到确定的回答。

% [原文 p.533]
\begin{theorem}[{见 {[}8{]}}]{2.1.3}
设 $X$ 是典范层（或反典范层）丰富的光滑射影簇。假设范畴 $\Db(X)$ 作为三角范畴等价于某个光滑代数簇 $X'$ 的导出范畴 $\Db(X')$。则 $X'$ 同构于 $X$。
\end{theorem}

[8] 中给出的这个定理的证明是构造性的，并给出了由凝聚层导出范畴还原簇的方法。此外，在定理的假设下，还可以只假设这两个导出范畴作为分次范畴等价，而不必作为三角范畴等价（见 [8]）。

在这种情形下，也可以描述正合自等价群。

\begin{theorem}[{见 {[}8{]}}]{2.1.4}
设 $X$ 是典范层（或反典范层）丰富的光滑射影簇。则范畴 $\Db(X)$ 的正合自等价的同构类构成的群，由簇的自同构、按线丛作扭转，以及导出范畴中的平移生成。
\end{theorem}

对任一簇 $X$，正合自等价构成的群 $\Auteq\Db(X)$ 总包含子群 $G(X)$，后者是正规子群 $G_1=\Pic(X)\oplus\Z$ 与自然作用于 $G_1$ 的子群 $G_2=\Aut X$ 的半直积。在这一包含 $G(X)\subset\Auteq\Db(X)$ 之下，$\Z$ 的生成元映到平移函子 $[1]$，线丛 $\cL\in\Pic(X)$ 映到函子 $\otimes\cL$，而自同构 $f:X\to X$ 诱导自等价 $\Rf f_*$。我们在 [8] 中证明：在定理 2.1.4 的假设下，正合自等价群 $\Auteq\Db(X)$ 等于 $G(X)$；也就是说，此时
\begin{equation*}
\Auteq\Db(X)\cong\Aut X\times(\Pic(X)\oplus\Z).
\end{equation*}

% [Tutor Note: 页图此处印作 $\Aut X\times(\Pic(X)\oplus\Z)$（乘号 $\times$），而文字层抽出的是半直积号 $\ltimes$。按上下文（$G(X)$ 为半直积）应以 $\ltimes$ 为是，属原书排版笔误；本译本照录页图的 $\times$，不静默改字。]

为研究两个簇何时具有等价的凝聚层导出范畴、并描述其自等价群，最好能有全部正合函子的显式公式。有一个猜想：它们总可由积上的对象代表，即总具有 (7) 的形式。在下一章中，我们对全忠实函子、从而对等价，给出这一猜想的证明。下一章整章都用于证明这一结果。于是，在研究光滑射影簇上凝聚层导出范畴之间的等价时，我们只需考虑形如 (7) 的函子。与解决这些问题相关的另一个问题，是需要一个判定给定函子是否为等价的准则。要证明函子 $F$ 是等价，只需证明 $F$ 与其右（或左）伴随都是全忠实函子（见定义 1.1.7）。

有一个办法可以判定函子 $\Phi_{\cE}:\Db(X)\to\Db(Y)$ 是否全忠实。

\begin{theorem}[{见 {[}7{]}}]{2.1.5}
设 $M$ 与 $X$ 是特征 $0$ 的代数闭域上的光滑射影簇，$\cE\in\Db(M\times X)$。此时函子 $\Phi_{\cE}$ 全忠实，当且仅当如下正交性条件成立：
\begin{enumerate}[label=\arabic*),leftmargin=2.4em,itemsep=1pt,topsep=2pt]
\item 对一切 $i$ 与一切 $t_1\neq t_2$，有 $\Hom^{i}_{X}(\Phi_{\cE}(\cO_{t_1}),\Phi_{\cE}(\cO_{t_2}))=0$；
% [原文 p.534]
\item {\raggedright $\Hom^{0}_{X}(\Phi_{\cE}(\cO_{t}),\Phi_{\cE}(\cO_{t}))=k$，且对任一 $i\notin\{0,\dim M\}$，有 $\Hom^{i}_{X}(\Phi_{\cE}(\cO_{t}),\Phi_{\cE}(\cO_{t}))=0$。\par}
\end{enumerate}
\end{theorem}

这里 $t$、$t_1$、$t_2$ 是 $M$ 的点，$\cO_{t_i}$ 是相应的摩天层。

这个定理的假设一般颇难验证；不过，当积上的对象 $\cE$ 是向量丛时，该准则相当好用。

考虑四个光滑完备代数簇 $X_1$、$X_2$、$Y_1$、$Y_2$。取分别属于范畴 $\Db(X_1\times Y_1)$ 与 $\Db(X_2\times Y_2)$ 的两个对象 $\cE_1$ 与 $\cE_2$，并考虑对象
\begin{equation*}
\cE_1\boxtimes\cE_2\in\Db\bigl((X_1\times X_2)\times(Y_1\times Y_2)\bigr),
\end{equation*}
按定义它就是 $p^{*}_{13}(\cE_1)\otL p^{*}_{24}(\cE_2)$。如上（见 (7)），对象 $\cE_1$、$\cE_2$ 与 $\cE_1\boxtimes\cE_2$ 定义函子
\begin{equation*}
\begin{gathered}
\Phi_{\cE_1}:\Db(X_1)\to\Db(Y_1),
\qquad
\Phi_{\cE_2}:\Db(X_2)\to\Db(Y_2),\\
\text{以及}\qquad
\Phi_{\cE_1\boxtimes\cE_2}:\Db(X_1\times X_2)\to\Db(Y_1\times Y_2).
\end{gathered}
\end{equation*}

我们考虑对象 $\cG\in\Db(X_1\times X_2)$，并用 $\cH$ 记对象 $\Phi_{\cE_1\boxtimes\cE_2}(\cG)\in\Db(Y_1\times Y_2)$。对这两个对象中的每一个，都可以按规则 (7) 配上函子：
\begin{equation*}
\Phi_{\cG}:\Db(X_1)\to\Db(X_2)
\qquad\text{与}\qquad
\Phi_{\cH}:\Db(Y_1)\to\Db(Y_2).
\end{equation*}

\begin{proposition}{2.1.6}
在上述记号下，有函子的同构
\begin{equation*}
\Phi_{\cH}\cong\Phi_{\cE_2}\circ\Phi_{\cG}\circ\Psi_{\cE_1}.
\end{equation*}
\end{proposition}

其证明立刻由命题 2.1.2 得到。

现在设 $Z_1$ 与 $Z_2$ 是另外两个光滑完备簇，而 $\cF_1$ 与 $\cF_2$ 分别是 $\Db(Y_1\times Z_1)$ 与 $\Db(Y_2\times Z_2)$ 中的对象；那么也有函子 $\Phi_{\cF_1}$、$\Phi_{\cF_2}$ 与 $\Phi_{\cF_1\boxtimes\cF_2}$。按规则 (9)，可以找到属于 $\Db(X_1\times Z_1)$ 与 $\Db(X_2\times Z_2)$ 的对象 $\cG_1$ 与 $\cG_2$，使得
\begin{equation*}
\Phi_{\cG_1}\cong\Phi_{\cF_1}\circ\Phi_{\cE_1}
\qquad\text{与}\qquad
\Phi_{\cG_2}\cong\Phi_{\cF_2}\circ\Phi_{\cE_2}.
\end{equation*}
直接验证表明有自然关系
\begin{equation}
\Phi_{\cF_1\boxtimes\cF_2}\circ\Phi_{\cE_1\boxtimes\cE_2}\cong\Phi_{\cG_1\boxtimes\cG_2}. \tag{10}
\end{equation}
利用这一点，容易证明下述断言。

\begin{proposition}{2.1.7}
在上述条件下，设 $\Phi_{\cE_1}$ 与 $\Phi_{\cE_2}$ 全忠实（相应地，是等价）。则函子
\begin{equation*}
\Phi_{\cE_1\boxtimes\cE_2}:\Db(X_1\times X_2)\to\Db(Y_1\times Y_2)
\end{equation*}
也全忠实（相应地，是范畴等价）。
\end{proposition}

% [原文 p.535]
\begin{proof}
若 $F$ 有左伴随 $F^{*}$（不妨这样说），则 $F$ 全忠实，当且仅当复合 $F^{*}F$ 同构于恒等函子。函子 $\Phi_{\cE_i}$ 有由 (8) 定义的左伴随 $\Phi^{*}_{\cE_i}$。既然它们全忠实，就推出复合 $\Phi^{*}_{\cE_i}\circ\Phi_{\cE_i}$ 同构于恒等函子，而恒等函子可由对角线 $\Delta_i\in X_i\times X_i$ 的结构层代表。容易看出，层 $\cO_{\Delta_1}\boxtimes\cO_{\Delta_2}$ 同构于对角线 $\Delta$ 的结构层 $\cO_{\Delta}$，其中 $\Delta$ 是 $(X_1\times X_2)\times(X_1\times X_2)$ 中的对角线。用公式 (10)，我们看到复合 $\Phi^{*}_{\cE_1\boxtimes\cE_2}\circ\Phi_{\cE_1\boxtimes\cE_2}$ 由对角线 $\Delta$ 的结构层代表，因而同构于恒等函子。于是 $\Phi_{\cE_1\boxtimes\cE_2}$ 全忠实。关于等价的断言可以类似地证明。
\end{proof}

现在假设函子 $\Phi_{\cE}:\Db(X)\to\Db(Y)$ 是等价，而 $\cF\in\Db(X\times Y)$ 是满足 $\Psi_{\cF}\cong\Phi_{\cE}^{-1}$ 的对象。由 (8)，有同构
\begin{equation*}
\cF\cong\cE^{\vee}\otimes p^{*}\omega_X[n]\cong\cE^{\vee}\otimes q^{*}\omega_Y[m],
\end{equation*}
它们立刻推出簇 $X$ 与 $Y$ 的维数 $n$ 与 $m$ 相等。

考虑函子
\begin{equation}
\Phi_{\cF\boxtimes\cE}:\Db(X\times X)\to\Db(Y\times Y) \tag{11}
\end{equation}
并把它记作 $Ad_{\cE}$。由命题 2.1.7，它也是等价。此外，由命题 2.1.6，对任一对象 $\cG\in\Db(X\times X)$ 有函子的同构
\begin{equation}
\Phi_{Ad_{\cE}(\cG)}\cong\Phi_{\cE}\circ\Phi_{\cG}\circ\Phi_{\cE}^{-1}. \tag{12}
\end{equation}

考虑 $\cG$ 是对角线的结构层 $\cO_{\Delta_X}$ 这一特殊情形，它代表恒等函子。于是，应用 (12)，我们看到函子 $Ad_{\cE}$ 把对角线的结构层 $\cO_{\Delta_X}$ 送到对角线的结构层 $\cO_{\Delta_Y}$。

考虑更一般的情形。用 $i_X$ 与 $i_Y$ 分别记对角线在 $X\times X$ 与 $Y\times Y$ 中的嵌入。我们把函子 $Ad_{\cE}$ 作用于对象 $i_{X*}\omega^{k}_{X}$，其中 $\omega_X$ 是 $X$ 的典范层（如上）。对象 $i_{X*}\omega^{k}_{X}$ 代表函子 $S^{k}[-nk]$，其中 $S$ 是 $\Db(X)$ 的 Serre 函子。既然每个等价都与 Serre 函子交换（引理 1.1.4），我们看到
\begin{equation}
Ad_{\cE}(i_{X*}\omega^{k}_{X})\cong i_{Y*}\omega^{k}_{Y}. \tag{13}
\end{equation}

\yizhu{$Ad_{\cE}$ 的做法是把等价 $\Phi_{\cE}:\Db(X)\to\Db(Y)$ 与它的逆夹在两边，从而把它“共轭”成 $\Db(X\times X)$ 与 $\Db(Y\times Y)$ 之间的等价：$\cG\mapsto\Phi_{\cE}\circ\Phi_{\cG}\circ\Phi_{\cE}^{-1}$，即 (12)。它与本节前半的复合律 (9) 是同一件事的两种说法——把 $\cF\boxtimes\cE$ 当核，就得到这个共轭。$Ad_{\cE}$ 保持对角线结构层（恒等函子的核），并把 $i_{X*}\omega^{k}_{X}$ 送到 $i_{Y*}\omega^{k}_{Y}$，这正是下面造 $\mathrm{HA}(X)$ 并证明它不变的依据。〔来源：原书 §2.1 式 (11)--(13)〕}

现在对每个簇 $X$ 定义双分次代数
\begin{equation*}
\mathrm{HA}(X)=\bigoplus_{i,k}\mathrm{HA}_{i,k}(X):=\bigoplus_{i,k}\Ext^{i}_{X\times X}(\cO_{\Delta_X},i_{X*}\omega^{k}_{X}).
\end{equation*}
这里的代数结构由 $\Ext$ 的复合给出，同时注意典范等同
\begin{equation*}
\Ext^{i}_{X\times X}(\cO_{\Delta_X},i_{X*}\omega^{k}_{X})\cong\Ext^{i}_{X\times X}(i_{X*}\omega^{m},i_{X*}\omega^{m+k}_{X}).
\end{equation*}

\yizhu{$\mathrm{HA}(X)$ 的分次指标有两个方向：$i$ 是 $\Ext$ 的次数，$k$ 是典范层的幂次。取 $i=0$ 的一行，按典范等同它给出 $\bigoplus_{k}H^{0}(X,\omega^{k}_{X})$，即 $X$ 的典范代数（$k\geqslant0$ 的部分）与反典范代数（$k\leqslant0$ 的部分）；这就是推论 2.1.9 说 $\mathrm{HA}$ 同时含典范代数与反典范代数的原因。原书引言里这个代数印作 $\mathrm{HA}^{i,k}$（上标），§2.1 此处及 (14) 印作 $\mathrm{HA}_{i,k}$（下标），本译本各照所据页图，未作统一。〔来源：原书引言 p.515 的 $\mathrm{HA}(X)$ 公式、§2.1 本处定义与推论 2.1.9〕}

为证明下一个定理，我们需要用到第 3 章的主要结果，它断言每个等价都可由积上的对象代表。

% [原文 p.536]
\begin{theorem}{2.1.8}
设 $X$ 与 $Y$ 是凝聚层导出范畴作为三角范畴等价的光滑射影簇。则双分次代数 $\mathrm{HA}(X)$ 与 $\mathrm{HA}(Y)$ 同构。
\end{theorem}

\begin{proof}
由定理 3.2.2，每个等价 $F:\Db(X)\to\Db(Y)$ 都可由积上的某个对象代表，因而同构于形如 $\Phi_{\cE}$ 的函子，其中 $\cE\in\Db(X\times Y)$。每个这样的等价定义等价
\begin{equation*}
Ad_{\cE}:\Db(X\times X)\to\Db(Y\times Y),
\end{equation*}
它把 $i_{X*}\omega^{k}_{X}$ 送到 $i_{Y*}\omega^{k}_{Y}$。等价 $Ad_{\cE}$ 诱导同构
\begin{equation*}
\Ext^{i}_{X\times X}(\cO_{\Delta_X},i_{X*}\omega^{k}_{X})\cong\Ext^{i}_{Y\times Y}(\cO_{\Delta_Y},i_{Y*}\omega^{k}_{Y}),
\end{equation*}
从而诱导双分次代数 $\mathrm{HA}(X)$ 与 $\mathrm{HA}(Y)$ 的同构。
\end{proof}

我们指出，由双分次代数 $\mathrm{HA}(X)$ 可以同时得到 $X$ 的典范代数与反典范代数。事实上，
\begin{equation*}
\bigoplus_{k\geqslant0}H^{0}(X,\omega^{k}_{X})=\bigoplus_{k\geqslant0}\mathrm{HA}_{0,k}(X)
\qquad\text{与}\qquad
\bigoplus_{k\leqslant0}H^{0}(X,\omega^{k}_{X})=\bigoplus_{k\leqslant0}\mathrm{HA}_{0,k}(X).
\end{equation*}
于是，定理 2.1.8 给出下述推论。

\begin{corollary}{2.1.9}
若两个光滑射影簇 $X$ 与 $Y$ 上凝聚层的导出范畴等价，则 $X$ 与 $Y$ 的典范代数（以及反典范代数）同构。
\end{corollary}

这个推论的陈述与定理 2.1.3 非常接近。不过我们应当指出，[8] 中给出的定理 2.1.3 的证明并不依赖下一章的主要结果，而且是构造性的。还要注意，在定理 2.1.3 中，我们并不假设第二个簇 $X'$ 的典范层（或反典范层）丰富；这一点由定理的证明推出。

我们也可以描述其余的空间 $\mathrm{HA}_{i,k}(X)$。在 [40] 中证明：计算
\begin{equation*}
\mathrm{HA}_{i,k}(X)=\Ext^{i}(\cO_{\Delta_X},i_{X*}\omega_X)
\end{equation*}
的谱序列，用 $\cO_{\Delta_X}$ 限制到对角线上的上同调表出，它在 $E_2$ 项退化。特别地，有同构
\begin{equation}
\mathrm{HA}_{i,k}(X)\cong\bigoplus_{p+q=i}H^{p}\Bigl(X,\bigwedge\nolimits^{q}T_X\otimes\omega^{k}_{X}\Bigr), \tag{14}
\end{equation}
其中 $T_X$ 是 $X$ 的切丛。此外，这个同构还变成代数的同构，也就是说，
\begin{equation*}
\mathrm{HA}(X)\cong\bigoplus_{i,k}\bigoplus_{p+q=i}H^{p}\Bigl(X,\bigwedge\nolimits^{q}T_X\otimes\omega^{k}_{X}\Bigr)
\end{equation*}
作为双分次代数同构。这一关系与定理 2.1.8 给出下述推论。

% [原文 p.537]
\begin{corollary}{2.1.10}
若两个光滑射影簇 $X$ 与 $Y$ 上凝聚层的导出范畴等价，则有向量空间的同构
\begin{equation}
\bigoplus_{p+q=i}H^{p}\Bigl(X,\bigwedge\nolimits^{q}T_X\otimes\omega^{k}\Bigr)
\cong
\bigoplus_{p+q=i}H^{p}\Bigl(Y,\bigwedge\nolimits^{q}T_Y\otimes\omega^{k}\Bigr). \tag{15}
\end{equation}
特别地，我们得到霍奇菱形各竖列之间的同构：
\begin{equation}
\bigoplus_{p-q=i}H^{p}(X,\Omega^{q}_{X})\cong\bigoplus_{p-q=i}H^{p}(Y,\Omega^{q}_{Y}). \tag{16}
\end{equation}
\end{corollary}

\begin{proof}
同构 (15) 立刻由定理 2.1.8 与等式 (14) 推出。同构 (16) 是 (15) 在 $k=1$ 时的特殊情形。
\end{proof}

霍奇菱形各竖列之间的同构也可以用另一种方式得到。设基域 $k$ 是 $\C$。对任一元素 $\xi\in H^{*}(X\times Y,\Q)$，我们可以定义线性映射
\begin{equation*}
v_{\xi}:H^{*}(X,\Q)\to H^{*}(Y,\Q)
\qquad\text{与}\qquad
w_{\xi}:H^{*}(Y,\Q)\to H^{*}(X,\Q),
\end{equation*}
其公式为
\begin{equation}
v_{\xi}(-)=q_{*}(\xi\cdot p^{*}(-))
\qquad\text{与}\qquad
w_{\xi}(-)=p_{*}(\xi\cdot q^{*}(-)). \tag{17}
\end{equation}

对这些映射，可以写出一个与函子复合的公式 (9) 类似的复合公式。设 $X$、$Y$、$Z$ 是三个光滑完备簇，而 $\xi$ 与 $\eta$ 分别是 $H^{*}(X\times Y,\Q)$ 与 $H^{*}(Y\times Z,\Q)$ 中的元素。则复合 $v_{\eta}\circ v_{\xi}$ 与映射 $v_{\zeta}$ 一致，其中 $\zeta\in H^{*}(X\times Z,\Q)$ 由公式
\begin{equation*}
\zeta=p_{XZ*}\bigl(p^{*}_{YZ}(\eta)\cup p^{*}_{XY}(\xi)\bigr)
\end{equation*}
给出。

对任一形如 $\Phi_{\cE}:\Db(X)\to\Db(Y)$ 的函子，我们可以配上线性映射 $\vp_{\cE}:H^{*}(X,\Q)\to H^{*}(Y,\Q)$。为此，按规则
\begin{equation}
\eps=p^{*}\sqrt{\td_X}\cdot\ch(\cE)\cdot q^{*}\sqrt{\td_Y} \tag{18}
\end{equation}
定义元素 $\eps\in H(X\times Y,\Q)$，其中 $\td_X$ 与 $\td_Y$ 分别是 $X$ 与 $Y$ 的 Todd 类，而 $\ch(\cE)$ 是 $\cE$ 的陈特征。我们定义映射
\begin{align}
\vp_{\cE}(-):=v_{\eps}(-)=q_{*}(\eps\cdot p^{*}(-)),\notag\\
\psi_{\cE}(-):=w_{\eps}(-)=p_{*}(\eps\cdot q^{*}(-)). \tag{19}
\end{align}

% [Tutor Note: 式 (18) 前后，页图与文字层都印作 $\eps\in H(X\times Y,\Q)$，缺了上标 $*$；按上下文应为 $H^{*}(X\times Y,\Q)$（(17) 处即作 $H^{*}$）。此处照录原样，不改字。]

\yizhu{$\vp_{\cE}$ 是函子 $\Phi_{\cE}$ 在“上同调层面”的影子：把核 $\cE$ 换成由 (18) 配出的类 $\eps$，再用 (17) 的 $v_{\eps}$ 作用。系数 $p^{*}\sqrt{\td_X}\cdot q^{*}\sqrt{\td_Y}$ 是为让复合保持复合（命题 2.1.11）而必需的修正项，其存在性由 Riemann--Roch 定理的 Grothendieck 形式保证；同一个修正项也保证推论 2.1.12 里的 K 群方块交换。$\psi_{\cE}$ 是反方向的对应物，只在需要沿 $p$ 推回时使用。〔来源：原书 §2.1 式 (18)--(19) 与命题 2.1.11 前的说明〕}

下一个命题立刻由 Riemann--Roch 定理的 Grothendieck 形式得到。

\begin{proposition}{2.1.11}
设 $\Phi_{\cE}:\Db(X)\to\Db(Z)$ 是某个复合 $\Phi_{\cG}\circ\Phi_{\cF}$，其中
\begin{equation*}
\Phi_{\cF}:\Db(X)\to\Db(Y),
\qquad
\Phi_{\cG}:\Db(Y)\to\Db(Z).
\end{equation*}
则 $\vp_{\cE}=\vp_{\cG}\circ\vp_{\cF}$。
\end{proposition}

这立刻给出下述推论。

% [原文 p.538]
\begin{corollary}{2.1.12}
若函子 $\Phi_{\cE}:\Db(X)\to\Db(Z)$ 是等价，则映射 $\vp_{\cE}:H^{*}(X,\Q)\to H^{*}(Y,\Q)$ 是同构，且其复化诱导霍奇菱形各竖列之间的同构 (16)。
\end{corollary}

% [Tutor Note: 推论 2.1.12 页图与文字层均作 $\Phi_{\cE}:\Db(X)\to\Db(Z)$，而同一句里的映射印作 $\vp_{\cE}:H^{*}(X,\Q)\to H^{*}(Y,\Q)$，$Z$ 与 $Y$ 不一致；据命题 2.1.11 的记号应为 $\Db(X)\to\Db(Y)$。此处照录原样。]

\begin{proof}
由命题 2.1.11，$\Phi_{\cE}$ 的拟逆函子诱导 $\vp_{\cE}$ 的逆映射。此外，由 (18)，元素 $\eps\in H^{*}(X\times Y,\Q)$ 对应一个代数闭链，因此容易验证 $\vp_{\cE}$ 的复化保持霍奇菱形的竖列。
\end{proof}

最后，我们还要指出，每个函子 $\Phi_{\cE}:\Db(X)\to\Db(Y)$ 都在范畴 $\Db(X)$ 与 $\Db(Y)$ 的 Grothendieck 群 $K(X)$ 与 $K(Y)$ 之间诱导映射 $\overline{\Phi}_{\cE}:K(X)\to K(Y)$。考虑把 $K(X)$ 的元素送到其陈特征乘以 Todd 类的平方根的映射
\begin{equation*}
\ch\cdot\sqrt{\td_X}:K(X)\to H^{*}(X,\Q),
\end{equation*}
利用 Riemann--Roch 定理可以证明，图表
\begin{center}
\begin{tikzcd}[row sep=3.0em,column sep=3.4em,ampersand replacement=\&]
K(X) \arrow[r, "\overline{\Phi}_{\cE}"] \arrow[d, "\ch\cdot\sqrt{\td_X}"] \& K(Y) \arrow[d, "\ch\cdot\sqrt{\td_Y}"] \\
H^{*}(X,\Q) \arrow[r, "\overline{\vp}_{\cE}"] \& H^{*}(Y,\Q)
\end{tikzcd}
\end{center}
交换。

% [Tutor Note: 本块（ch02a）按分块约定止于 “2.2. Examples of equivalences: flopping birational transformations” 标题之前；该标题及其后内容属 ch02b。推论 2.1.12 与该 K 群交换图排在原书 p.538，故本块末页为 p.538（页数 9 页，而非分块表的 8 页），特此登记，供 Lead 校边界。]
